\begin{table}[h]
\centering
\caption{Cálculo de EIAC y clasificación funcional de columnas}
\label{tab:calculo-eiac-columnas-20260624_095401}
\scriptsize
\begin{tabular}{|p{2.1cm}|p{3.0cm}|c|c|c|c|c|c|p{2.5cm}|p{2.8cm}|}
\hline
\textbf{Columna} & \textbf{Distribución} & \textbf{\(k\)} & \textbf{\(H(X)\)} & \textbf{\(H_n(X)\)} & \textbf{\(C(X)\)} & \textbf{\(U(X)\)} & \textbf{EIAC} & \textbf{Análisis} & \textbf{Clasificación funcional} \\
\hline
Abstraction & Detailed = 341; Standard = 197; Meta = 77 & 3 & 1.373 & 0.866 & 1.000 & 0.005 & 0.862 & Alta relevancia & Analítica \\
\hline
consequences\_scope\_count & 0 = 246; 1 = 109; 4 = 73; 3 = 45; ... + 8 valores & 12 & 2.724 & 0.760 & 1.000 & 0.020 & 0.745 & Alta relevancia & Analítica complementaria \\
\hline
consequences\_impact\_count & 0 = 247; 1 = 140; 3 = 96; 2 = 53; ... + 5 valores & 9 & 2.297 & 0.725 & 1.000 & 0.015 & 0.714 & Alta relevancia & Analítica complementaria \\
\hline
mitigations\_count & 0 = 215; 1 = 131; 2 = 96; 3 = 62; ... + 9 valores & 13 & 2.675 & 0.723 & 1.000 & 0.021 & 0.708 & Alta relevancia & Analítica complementaria \\
\hline
skills\_required\_count & 0 = 319; 1 = 207; 2 = 72; 3 = 16; ... + 1 valores & 5 & 1.534 & 0.661 & 1.000 & 0.008 & 0.655 & Relevancia media & Analítica complementaria \\
\hline
prerequisites\_count & 1 = 330; 0 = 125; 2 = 95; 3 = 39; ... + 3 valores & 7 & 1.846 & 0.658 & 1.000 & 0.011 & 0.650 & Relevancia media & Analítica complementaria \\
\hline
Typical Severity & High = 211; Medium = 112; Low = 104; Very High = 59; ... + 1 valores & 5 & 1.909 & 0.822 & 0.797 & 0.008 & 0.650 & Relevancia media & Analítica \\
\hline
related\_weaknesses\_count & 1 = 255; 0 = 165; 2 = 56; 3 = 37; ... + 12 valores & 16 & 2.544 & 0.636 & 1.000 & 0.026 & 0.619 & Relevancia media & Analítica complementaria \\
\hline
example\_instances\_count & 0 = 346; 1 = 174; 2 = 76; 3 = 12; ... + 2 valores & 6 & 1.546 & 0.598 & 1.000 & 0.010 & 0.592 & Relevancia media & Analítica complementaria \\
\hline
resources\_required\_count & 0 = 343; 1 = 258; 2 = 10; 3 = 4 & 4 & 1.139 & 0.570 & 1.000 & 0.007 & 0.566 & Relevancia media & Analítica complementaria \\
\hline
execution\_flow\_steps\_count & 0 = 376; 3 = 102; 4 = 85; 2 = 30; ... + 4 valores & 8 & 1.703 & 0.568 & 1.000 & 0.013 & 0.560 & Relevancia media & Analítica complementaria \\
\hline
Likelihood Of Attack & High = 131; Medium = 118; Low = 96 & 3 & 1.573 & 0.993 & 0.561 & 0.005 & 0.554 & Relevancia media & Analítica \\
\hline
Status & Draft = 402; Stable = 154; Deprecated = 56; Usable = 2; ... + 1 valores & 5 & 1.258 & 0.542 & 1.000 & 0.008 & 0.537 & Relevancia media & Analítica \\
\hline
related\_attack\_patterns\_count & 1 = 398; 0 = 109; 2 = 58; 3 = 27; ... + 5 valores & 9 & 1.625 & 0.513 & 1.000 & 0.015 & 0.505 & Relevancia media & Analítica complementaria \\
\hline
Related Attack Patterns & ::NATURE:ChildOf:CAPEC ID:312:: = 16; ::NATURE:ChildOf:CAPEC ID:444:: = 14; ::NATURE:ChildOf:CAPEC ID:169:: = 12; ::NATURE:ChildOf:CAPEC ID:267:: = 11; ... + 208 valores & 212 & 7.236 & 0.936 & 0.823 & 0.345 & 0.505 & Relevancia media & Analítica complementaria \\
\hline
execution\_flow\_techniques\_count & 0 = 414; 3 = 30; 5 = 26; 6 = 25; ... + 13 valores & 17 & 2.059 & 0.504 & 1.000 & 0.028 & 0.490 & Relevancia media & Analítica complementaria \\
\hline
taxonomy\_mappings\_count & 0 = 392; 1 = 149; 2 = 46; 3 = 16; ... + 4 valores & 8 & 1.471 & 0.490 & 1.000 & 0.013 & 0.484 & Relevancia media & Analítica complementaria \\
\hline
Related Weaknesses & ::200:: = 51; ::284:: = 15; ::770:: = 10; ::829:: = 9; ... + 266 valores & 270 & 7.344 & 0.909 & 0.732 & 0.439 & 0.373 & Relevancia media & Analítica complementaria \\
\hline
Consequences & ::SCOPE:Confidentiality:TECHNICAL IMPACT:Read Data:: = 25; ::SCOPE:Confidentiality:TECHNICAL IMPACT:Other::SCOPE:Confidentiality:SCOPE:Access Control:SCOPE:Authorization:TECHNICAL IMPACT:Bypass Protection Mechanism:TECHNICAL IMPACT:Hide Activities:: = 18; ::SCOPE:Confidentiality:TECHNICAL IMPACT:Read Data::SCOPE:Confidentiality:SCOPE:Access Control:SCOPE:Authorization:TECHNICAL IMPACT:Bypass Protection Mechanism:TECHNICAL IMPACT:Hide Activities:: = 15; ::SCOPE:Confidentiality:SCOPE:Access Control:SCOPE:Authorization:TECHNICAL IMPACT:Gain Privileges:: = 11; ... + 216 valores & 220 & 7.168 & 0.921 & 0.600 & 0.358 & 0.355 & Relevancia media & Analítica complementaria \\
\hline
description\_length\_chars & 156 = 7; 174 = 6; 258 = 6; 40 = 6; ... + 419 valores & 423 & 8.536 & 0.978 & 1.000 & 0.688 & 0.305 & Baja relevancia & Candidata a limpieza o revisión \\
\hline
Skills Required & ::SKILL:The adversary requires strong inter-personal and communication skills.:LEVEL:Low:: = 11; ::SKILL:Detailed knowledge on HTTP protocol: request and response messages structure and usage of specific headers.:LEVEL:Medium::SKILL:Detailed knowledge on how specific HTTP agents receive, send, process, interpret, and parse a variety of HTTP messages and headers.:LEVEL:Medium::SKILL:Possess knowledge on the exact details in the discrepancies between several targeted HTTP agents in path of an HTTP message in parsing its message structure and individual headers.:LEVEL:Medium:: = 4; ::SKILL:To identify and execute against an over-privileged system interface:LEVEL:Low:: = 4; ::SKILL:An adversary can simply overflow a buffer by inserting a long string into an adversary-modifiable injection vector. The result can be a DoS.:LEVEL:Low::SKILL:Exploiting a buffer overflow to inject malicious code into the stack of a software system or even the heap can require a higher skill level.:LEVEL:High:: = 4; ... + 258 valores & 262 & 7.889 & 0.982 & 0.481 & 0.426 & 0.271 & Baja relevancia & Bajo aporte por datos faltantes \\
\hline
Resources Required & ::None: No specialized resources are required to execute this type of attack.:: = 54; ::A tool capable of sending and receiving packets from a remote system.:: = 7; ::Scanners or utilities that provide the ability to send custom ICMP queries.:: = 3; ::Any type of active probing that involves non-standard packet headers requires the use of raw sockets, which is not available on particular operating systems (Microsoft Windows XP SP 2, for example). Raw socket manipulation on Unix/Linux requires root privileges. A tool capable of sending and receiving packets from a remote system.:: = 3; ... + 196 valores & 200 & 6.772 & 0.886 & 0.442 & 0.325 & 0.264 & Baja relevancia & Bajo aporte por datos faltantes \\
\hline
Mitigations & ::An organization should provide regular, robust cybersecurity training to its employees to prevent social engineering attacks.:: = 8; ::Design: evaluate HTTP agents prior to deployment for parsing/interpretation discrepancies.::Configuration: front-end HTTP agents notice ambiguous requests.::Configuration: back-end HTTP agents reject ambiguous requests and close the network connection.::Configuration: Disable reuse of back-end connections.::Configuration: Use HTTP/2 for back-end connections.::Configuration: Use the same web server software for front-end and back-end server.::Implementation: Utilize a Web Application Firewall (WAF) that has built-in mitigation to detect abnormal requests/responses.::Configuration: Install latest vendor security patches available for both intermediary and back-end HTTP infrastructure (i.e. proxies and web servers)::Configuration: Ensure that HTTP infrastructure in the chain or network path utilize a strict uniform parsing process.::Implementation: Utilize intermediary HTTP infrastructure capable of filtering and/or sanitizing user-input.:: = 2; ::Assume all input is malicious. Create an allowlist that defines all valid input to the software system based on the requirements specifications. Input that does not match against the allowlist should not be permitted to enter into the system. Test your decoding process against malicious input.::Be aware of the threat of alternative method of data encoding and obfuscation technique such as IP address encoding.::When client input is required from web-based forms, avoid using the GET method to submit data, as the method causes the form data to be appended to the URL and is easily manipulated. Instead, use the POST method whenever possible.::Any security checks should occur after the data has been decoded and validated as correct data format. Do not repeat decoding process, if bad character are left after decoding process, treat the data as suspicious, and fail the validation process.::Refer to the RFCs to safely decode URL.::Regular expression can be used to match safe URL patterns. However, that may discard valid URL requests if the regular expression is too restrictive.::There are tools to scan HTTP requests to the server for valid URL such as URLScan from Microsoft (http://www.microsoft.com/technet/security/tools/urlscan.mspx).:: = 2; ::Design: Build throttling mechanism into the resource allocation. Provide for a timeout mechanism for allocated resources whose transaction does not complete within a specified interval.::Implementation: Provide for network flow control and traffic shaping to control access to the resources.:: = 2; ... + 372 valores & 376 & 8.499 & 0.993 & 0.650 & 0.611 & 0.251 & Baja relevancia & Candidata a limpieza o revisión \\
\hline
Example Instances & ::The attacker uses relative path traversal to access files in the application. This is an example of accessing user's password file. http://www.example.com/getProfile.jsp?filename=../../../../etc/passwd However, the target application employs regular expressions to make sure no relative path sequences are being passed through the application to the web page. The application would replace all matches from this regex with the empty string. Then an attacker creates special payloads to bypass this filter: http://www.example.com/getProfile.jsp?filename=\%2e\%2e/\%2e\%2e/\%2e\%2e/\%2e\%2e /etc/passwd When the application gets this input string, it will be the desired vector by the attacker.:: = 2; ::Consider the case of attack performed against the createCustomerBillingAccount Web Service for an online store. In this case, the createCustomerBillingAccount Web Service receives a huge number of simultaneous requests, containing nonsense billing account creation information (the small XML messages). The createCustomerBillingAccount Web Services may forward the messages to other Web Services for processing. The application suffers from a high load of requests, potentially leading to a complete loss of availability the involved Web Service.:: = 2; ::Consider an application that uses an XML database to authenticate its users. The application retrieves the user name and password from a request and forms an XPath expression to query the database. An attacker can successfully bypass authentication and login without valid credentials through XPath Injection. This can be achieved by injecting the query to the XML database with XPath syntax that causes the authentication check to fail. Improper validation of user-controllable input and use of a non-parameterized XPath expression enable the attacker to inject an XPath expression that causes authentication bypass.:: = 2; ::Implementing the Model-View-Controller (MVC) within Java EE's Servlet paradigm using a Single front controller pattern that demands that brokered HTTP requests be authenticated before hand-offs to other Action Servlets. If no security-constraint is placed on those Action Servlets, such that positively no one can access them, the front controller can be subverted.:: = 1; ... + 262 valores & 266 & 8.049 & 0.999 & 0.437 & 0.433 & 0.248 & Baja relevancia & Bajo aporte por datos faltantes \\
\hline
indicators\_count & 0 = 559; 1 = 36; 3 = 8; 2 = 8; ... + 1 valores & 5 & 0.575 & 0.248 & 1.000 & 0.008 & 0.246 & Baja relevancia & Bajo aporte por baja variabilidad \\
\hline
Taxonomy Mappings & TAXONOMY NAME:ATTACK:ENTRY ID:1195.003:ENTRY NAME:Supply Chain Compromise: Compromise Hardware Supply Chain:: = 10; TAXONOMY NAME:ATTACK:ENTRY ID:1195.002:ENTRY NAME:Supply Chain Compromise: Compromise Software Supply Chain:: = 5; TAXONOMY NAME:ATTACK:ENTRY ID:1548:ENTRY NAME:Abuse Elevation Control Mechanism:: = 4; TAXONOMY NAME:ATTACK:ENTRY ID:1082:ENTRY NAME:System Information Discovery:: = 3; ... + 187 valores & 191 & 7.410 & 0.978 & 0.363 & 0.311 & 0.244 & Baja relevancia & Bajo aporte por datos faltantes \\
\hline
Execution Flow & ::STEP:1:PHASE:Explore:DESCRIPTION:[Survey] The attacker surveys the target application, possibly as a valid and authenticated user:TECHNIQUE:Spidering web sites for all available links:TECHNIQUE:Brute force guessing of resource names:TECHNIQUE:Brute force guessing of user names / credentials:TECHNIQUE:Brute force guessing of function names / actions::STEP:2:PHASE:Explore:DESCRIPTION:[Identify Functionality] At each step, the attacker notes the resource or functionality access mechanism invoked upon performing specific actions:TECHNIQUE:Use the web inventory of all forms and inputs and apply attack data to those inputs.:TECHNIQUE:Use a packet sniffer to capture and record network traffic:TECHNIQUE:Execute the software in a debugger and record API calls into the operating system or important libraries. This might occur in an environment other than a production environment, in order to find weaknesses that can be exploited in a production environment.::STEP:3:PHASE:Experiment:DESCRIPTION:[Iterate over access capabilities] Possibly as a valid user, the attacker then tries to access each of the noted access mechanisms directly in order to perform functions not constrained by the ACLs.:TECHNIQUE:Fuzzing of API parameters (URL parameters, OS API parameters, protocol parameters):: = 1; ::STEP:1:PHASE:Explore:DESCRIPTION:[Identify target application] The adversary identifies a target application or program to perform the buffer overflow on. In this attack the adversary looks for an application that loads the content of an environment variable into a buffer.::STEP:2:PHASE:Experiment:DESCRIPTION:[Find injection vector] The adversary identifies an injection vector to deliver the excessive content to the targeted application's buffer.:TECHNIQUE:Change the values of environment variables thought to be used by the application to contain excessive data. If the program is loading the value of the environment variable into a buffer, this could cause a crash and an attack vector will be found.::STEP:3:PHASE:Experiment:DESCRIPTION:[Craft overflow content] The adversary crafts the content to be injected. If the intent is to simply cause the software to crash, the content need only consist of an excessive quantity of random data. If the intent is to leverage the overflow for execution of arbitrary code, the adversary crafts the payload in such a way that the overwritten return address is replaced with one of the adversary's choosing.:TECHNIQUE:Create malicious shellcode that will execute when the program execution is returned to it.:TECHNIQUE:Use a NOP-sled in the overflow content to more easily slide into the malicious code. This is done so that the exact return address need not be correct, only in the range of all of the NOPs::STEP:4:PHASE:Exploit:DESCRIPTION:[Overflow the buffer] Using the injection vector, the adversary injects the crafted overflow content into the buffer.:: = 1; ::STEP:1:PHASE:Explore:DESCRIPTION:[Identify target application] The adversary identifies a target application or program to perform the buffer overflow on. Adversaries often look for applications that accept user input and that perform manual memory management.::STEP:2:PHASE:Experiment:DESCRIPTION:[Find injection vector] The adversary identifies an injection vector to deliver the excessive content to the targeted application's buffer.:TECHNIQUE:Provide large input to a program or application and observe the behavior. If there is a crash, this means that a buffer overflow attack is possible.::STEP:3:PHASE:Experiment:DESCRIPTION:[Craft overflow content] The adversary crafts the content to be injected. If the intent is to simply cause the software to crash, the content need only consist of an excessive quantity of random data. If the intent is to leverage the overflow for execution of arbitrary code, the adversary crafts the payload in such a way that the overwritten return address is replaced with one of the adversary's choosing.:TECHNIQUE:Create malicious shellcode that will execute when the program execution is returned to it.:TECHNIQUE:Use a NOP-sled in the overflow content to more easily slide into the malicious code. This is done so that the exact return address need not be correct, only in the range of all of the NOPs::STEP:4:PHASE:Exploit:DESCRIPTION:[Overflow the buffer] Using the injection vector, the adversary injects the crafted overflow content into the buffer.:: = 1; ::STEP:1:PHASE:Explore:DESCRIPTION:[Determine applicability] The adversary determines whether server side includes are enabled on the target web server.:TECHNIQUE:Look for popular page file names. The attacker will look for .shtml, .shtm, .asp, .aspx, and other well-known strings in URLs to help determine whether SSI functionality is enabled.:TECHNIQUE:Fetch .htaccess file. In Apache web server installations, the .htaccess file may enable server side includes in specific locations. In those cases, the .htaccess file lives inside the directory where SSI is enabled, and is theoretically fetchable from the web server. Although most web servers deny fetching the .htaccess file, a misconfigured server will allow it. Thus, an attacker will frequently try it.::STEP:2:PHASE:Experiment:DESCRIPTION:[Find Injection Point] Look for user controllable input, including HTTP headers, that can carry server side include directives to the web server.:TECHNIQUE:Use a spidering tool to follow and record all links. Make special note of any links that include parameters in the URL.:TECHNIQUE:Use a proxy tool to record all links visited during a manual traversal of the web application. Make special note of any links that include parameters in the URL. Manual traversal of this type is frequently necessary to identify forms that are GET method forms rather than POST forms.::STEP:3:PHASE:Exploit:DESCRIPTION:[Inject SSI] Using the found injection point, the adversary sends arbitrary code to be inlcuded by the application on the server side. They may then need to view a particular page in order to have the server execute the include directive and run a command or open a file on behalf of the adversary.:: = 1; ... + 235 valores & 239 & 7.901 & 1.000 & 0.389 & 0.389 & 0.238 & Baja relevancia & Bajo aporte por datos faltantes \\
\hline
Prerequisites & ::None:: = 11; ::The ability to monitor and interact with network communications.Access to at least one host, and the privileges to interface with the network interface card.:: = 11; ::The adversary must have the means and knowledge of how to communicate with the target in some manner.:: = 9; ::Targeted software is utilizing application framework APIs:: = 6; ... + 431 valores & 435 & 8.590 & 0.980 & 0.797 & 0.707 & 0.229 & Baja relevancia & Candidata a limpieza o revisión \\
\hline
alternate\_terms\_count & 0 = 598; 2 = 9; 1 = 3; 3 = 2; ... + 2 valores & 6 & 0.235 & 0.091 & 1.000 & 0.010 & 0.090 & Baja relevancia & Bajo aporte por baja variabilidad \\
\hline
Indicators & ::Differences in requests processed by the two agents. This requires careful monitoring or a capable log analysis tool.:: = 2; ::Many incorrect login attempts are detected by the system.:: = 2; ::Authentication attempts use credentials that have been used previously by the account in question.::Authentication attempts are originating from IP addresses or locations that are inconsistent with the user's normal IP addresses or locations.::Data is being transferred and/or removed from systems/applications within the network.::Suspicious or Malicious software is downloaded/installed on systems within the domain.::Messages from a legitimate user appear to contain suspicious links or communications not consistent with the user's normal behavior.:: = 2; ::Many invalid login attempts are coming from the same machine (same IP address) or for multiple user accounts within short succession.::The login attempts use passwords that have been used previously by the user account in question.::Login attempts are originating from IP addresses or locations that are inconsistent with the user's normal IP addresses or locations.:: = 2; ... + 46 valores & 50 & 5.593 & 0.991 & 0.091 & 0.081 & 0.083 & Baja relevancia & Bajo aporte por datos faltantes \\
\hline
notes\_length\_chars & 0 = 596; 718 = 4; 509 = 2; 602 = 1; ... + 12 valores & 16 & 0.314 & 0.078 & 1.000 & 0.026 & 0.076 & Baja relevancia & Bajo aporte por baja variabilidad \\
\hline
Description & This attack pattern has been deprecated. = 6; This attack pattern has been deprecated as it was deemed not to be a legitimate attack pattern. Please refer to CAPEC-118 : Collect and Analyze Information. = 4; This attack pattern has been deprecated as it is a duplicate of the existing attack pattern CAPEC-65 : Sniff Application Code. Please refer to this other CAPEC going forward. = 3; This CAPEC has been deprecated because it is not directly related to a weakness, social engineering, supply chains, or a physical-based attack. = 3; ... + 591 valores & 595 & 9.186 & 0.997 & 0.997 & 0.967 & 0.032 & Baja relevancia & Descriptiva contextual \\
\hline
Notes & TYPE:Terminology:NOTE:HTTP Splitting – the act of forcing a sender of (HTTP) messages to emit data stream consisting of more messages than the sender’s intension. The messages sent are 100\% valid and RFC compliant [REF-117].::::TYPE:Terminology:NOTE:HTTP Smuggling – the act of forcing a sender of (HTTP) messages to emit data stream which may be parsed as a different set of messages (i.e. dislocated message boundaries) than the sender’s intention. This is done by virtue of forcing the sender to emit non-standard messages which can be interpreted in more than one way [REF-117].::::TYPE:Relationship:NOTE:HTTP Smuggling is an evolution of previous HTTP Splitting techniques which are commonly remediated against.:: = 4; TYPE:Other:NOTE:Large quantities of data is often moved from the target system to some other adversary controlled system. Data found on a target system might require extensive resources to be fully analyzed. Using these resources on the target system might enable a defender to detect the adversary. Additionally, proper analysis tools required might not be available on the target system.::::TYPE:Other:NOTE:This attack differs from Data Interception and other data collection attacks in that the attacker actively queries the target rather than simply watching for the target to reveal information.:: = 1; TYPE:Other:NOTE:Other class of attacks focus on firmware, where malicious updates are made to the core system firmware or BIOS. Since this occurs outside the controls of the operating system, the OS detection and prevention mechanisms do not aid, thus allowing an adversary to evade defenses as well as gain persistence on the target's system.:: = 1; TYPE:Other:NOTE:White box analysis techniques include file or binary analysis, debugging, disassembly, and decompilation, and generally fall into categories referred to as 'static' and 'dynamic' analysis. Static analysis encompasses methods which analyze the binary, or extract its source code or object code without executing the program. Dynamic analysis involves analyzing the program during execution. Some forms of file analysis tools allow the executable itself to be analyzed, the most basic of which can analyze features of the binary. More sophisticated forms of static analysis analyze the binary file and extract assembly code, and possibly source code representations, from analyzing the structure of the file itself. Dynamic analysis tools execute the binary file and monitor its in memory footprint, revealing its execution flow, memory usage, register values, and machine instructions. This type of analysis is most effective for analyzing the execution of binary files whose content has been obfuscated or encrypted in its native executable form. Debuggers allow the program's execution to be monitored, and depending upon the debugger's sophistication may show relevant source code for each step in execution, or may display and allow interactions with memory, variables, or values generated by the program during run-time operations. Disassemblers operate in reverse of assemblers, allowing assembly code to be extracted from a program as it executes machine code instructions. Disassemblers allow low-level interactions with the program as it executes, such as manipulating the program's run time operations. Decompilers can be utilized to analyze a binary file and extract source code from the compiled executable. Collectively, the tools and methods described are those commonly applied to a binary executable file and provide means for reverse engineering the file by revealing the hidden functions of its operation or composition.:: = 1; ... + 12 valores & 16 & 3.827 & 0.957 & 0.031 & 0.026 & 0.029 & Baja relevancia & Descriptiva contextual \\
\hline
Alternate Terms & ::TERM:XML Denial of Service (XML DoS):DESCRIPTION::: = 3; ::TERM:HTTP Desync:DESCRIPTION:Modification/manipulation of HTTP message headers, request-line and body parameters to disrupt and interfere in the interpretation and parsing of HTTP message lengths/boundaries for consecutive HTTP messages by HTTP agents in a HTTP chain or network path.:: = 2; ::TERM:Directory Traversal:DESCRIPTION::: = 1; ::TERM:Smishing:DESCRIPTION:::TERM:MobPhishing:DESCRIPTION::: = 1; ... + 10 valores & 14 & 3.690 & 0.969 & 0.028 & 0.023 & 0.026 & Baja relevancia & Descriptiva contextual \\
\hline
ID & 1 = 1; 10 = 1; 100 = 1; 101 = 1; ... + 611 valores & 615 & 9.264 & 1.000 & 1.000 & 1.000 & 0.000 & Baja relevancia & Trazabilidad \\
\hline
Name & Accessing Functionality Not Properly Constrained by ACLs = 1; Buffer Overflow via Environment Variables = 1; Overflow Buffers = 1; Server Side Include (SSI) Injection = 1; ... + 611 valores & 615 & 9.264 & 1.000 & 1.000 & 1.000 & 0.000 & Baja relevancia & Trazabilidad \\
\hline
\end{tabular}
\end{table}
